\documentclass[tikz,border=8pt]{standalone}
\usepackage{tikz}
\usepackage{amsmath}
\usepackage{amssymb}
\usepackage{ctex}
\usetikzlibrary{calc, positioning, arrows.meta, decorations.pathreplacing}

\begin{document}
\begin{tikzpicture}[>=Stealth]

%% ===== 坐标系 =====
\draw[gray!25, thin] (-0.5,-1.5) grid (4.5,3.5);
\draw[->,gray!60, thin] (-0.6,0) -- (4.8,0) node[right,font=\scriptsize] {$x$};
\draw[->,gray!60, thin] (0,-1.6) -- (0,3.7) node[above,font=\scriptsize] {$y$};
\fill (0,0) circle (1.5pt) node[below left,font=\scriptsize] {$O$};

%% ===== 蓝色标准基 =====
\draw[->,thick,blue!80!black] (0,0) -- (1,0)
  node[below,font=\small,blue!80!black] {$\mathbf{e}_1$};
\draw[->,thick,blue!60] (0,0) -- (0,1)
  node[left,font=\small,blue!60] {$\mathbf{e}_2$};
\node[font=\scriptsize,blue!80!black] at (1.0,-0.9) {标准基 $\mathcal{B}$};

%% ===== 红色斜坐标基 b1=(1,1), b2=(1,-1) =====
\draw[->,thick,red!80!black] (0,0) -- (1,1)
  node[right,font=\small,red!80!black] {$\mathbf{b}_1=(1,1)$};
\draw[->,thick,red!60] (0,0) -- (1,-1)
  node[right,font=\small,red!60] {$\mathbf{b}_2=(1,-1)$};
\node[font=\scriptsize,red!80!black] at (2.2,-1.3) {斜坐标基 $\mathcal{B}'$};

%% ===== 向量 v =====
% 取 v = (2, 1)
% 在标准基下：[v]_B = (2,1)
% 在新基下：v = a*(1,1) + b*(1,-1) => a+b=2, a-b=1 => a=1.5, b=0.5
% [v]_{B'} = (1.5, 0.5)

\draw[->,very thick,purple!80!black] (0,0) -- (2,1)
  node[above right,font=\small,purple!80!black] {$\mathbf{v}=(2,1)$};
\fill[purple!70] (2,1) circle (2pt);

%% ===== 标准基坐标分解（蓝色虚线）=====
\draw[blue!40, dashed, thin] (2,1) -- (2,0) node[below,font=\scriptsize,blue!70] {$2$};
\draw[blue!40, dashed, thin] (2,1) -- (0,1) node[left,font=\scriptsize,blue!70] {$1$};

%% ===== 斜坐标分解（红色虚线）：1.5*(1,1) + 0.5*(1,-1) =====
\draw[red!40, dashed, thin] (0,0) -- (1.5,1.5);
\draw[red!40, dashed, thin] (1.5,1.5) -- (2,1);
% 标注分量
\node[font=\scriptsize,red!70] at (0.55,1.0) {$1.5\,\mathbf{b}_1$};
\node[font=\scriptsize,red!70] at (2.1,1.4) {$0.5\,\mathbf{b}_2$};

%% ===== 坐标标注框 =====
\node[fill=blue!8, draw=blue!40, rounded corners=4pt, inner sep=6pt,
      align=center, font=\small] at (4.0, 2.8) {
  $[\mathbf{v}]_{\mathcal{B}} = \begin{pmatrix}2\\1\end{pmatrix}$
};

\node[fill=red!8, draw=red!40, rounded corners=4pt, inner sep=6pt,
      align=center, font=\small] at (4.0, 1.5) {
  $[\mathbf{v}]_{\mathcal{B}'} = \begin{pmatrix}3/2\\1/2\end{pmatrix}$
};

%% ===== 过渡矩阵公式框（外侧）=====
\node[font=\small, fill=white, draw=gray!50, thin, rounded corners=4pt,
      inner sep=8pt, align=center]
  at (2.0, -2.8) {
  \textbf{过渡矩阵 $P$：} \quad $P = \begin{pmatrix}1&1\\1&-1\end{pmatrix}$（列 = 新基向量）\\[4pt]
  $[\mathbf{v}]_{\mathcal{B}'} = P^{-1}[\mathbf{v}]_{\mathcal{B}}$
  \qquad
  $P^{-1} = \dfrac{1}{2}\begin{pmatrix}1&1\\1&-1\end{pmatrix}$
};

%% ===== 整体标题 =====
\node[font=\large\bfseries] at (2.0, 4.4) {基变换与坐标变化};

\end{tikzpicture}
\end{document}
